
\documentclass[10pt]{article}
\usepackage[utf8]{inputenc}
\usepackage[T1]{fontenc}
\usepackage{amsmath, amssymb, xcolor, geometry, enumitem, multicol, tcolorbox}

\definecolor{mainblue}{RGB}{0, 102, 204}
\geometry{a4paper, margin=1.2cm, top=1cm, bottom=3cm, footskip=1.2cm}

\begin{document}
\enlargethispage{2.5cm}

% Modern Header
\begin{tcolorbox}[colback=mainblue!5, colframe=mainblue, sharp corners, boxrule=0.5pt]
    \small \textbf{Unit:} undefined \hfill \textbf{Name:} \rule{3cm}{0.4pt} \\
    \small \textbf{Topic:} Variables and Expressions / Order of Operations \hfill \textbf{Date:} \rule{2cm}{0.4pt} \\
    \centering \textbf{\large Challenge (Level 0)}
\end{tcolorbox}

\vspace{0.2cm}

% MCQ Section
\begin{multicols}{2}
\begin{enumerate}[label=\textbf{Q\arabic*.}, leftmargin=*, itemsep=6pt, parsep=0pt]

  \item \small What is the variable in the expression $2x + 3$?
  \begin{enumerate}[label=(\Alph*), leftmargin=0.6cm, nosep, itemsep=2pt]
    \item 2
    \item 3
    \item x
    \item 2x
    \item 3x
  \end{enumerate}

  \item \small Identify the constant term in the expression $4y - 2$.
  \begin{enumerate}[label=(\Alph*), leftmargin=0.6cm, nosep, itemsep=2pt]
    \item 4y
    \item 4
    \item -2
    \item y
    \item 2y
  \end{enumerate}

  \item \small What is the coefficient of $x$ in the expression $5x - 1$?
  \begin{enumerate}[label=(\Alph*), leftmargin=0.6cm, nosep, itemsep=2pt]
    \item 1
    \item -1
    \item 5
    \item x
    \item -5
  \end{enumerate}

  \item \small Which of the following is an example of a simple expression?
  \begin{enumerate}[label=(\Alph*), leftmargin=0.6cm, nosep, itemsep=2pt]
    \item $2x + 3 = 5$
    \item $4y - 2$
    \item $x^2 + 3x - 2$
    \item $rac{2}{3}x + 1$
    \item $x = 2$
  \end{enumerate}

  \item \small What operation is represented by the symbol $+$?
  \begin{enumerate}[label=(\Alph*), leftmargin=0.6cm, nosep, itemsep=2pt]
    \item Subtraction
    \item Addition
    \item Multiplication
    \item Division
    \item Exponentiation
  \end{enumerate}

  \item \small Identify the operator in the expression $7 - 2$.
  \begin{enumerate}[label=(\Alph*), leftmargin=0.6cm, nosep, itemsep=2pt]
    \item 7
    \item 2
    \item -
    \item +
    \item *
  \end{enumerate}

  \item \small What is the term $3x$ an example of in an algebraic expression?
  \begin{enumerate}[label=(\Alph*), leftmargin=0.6cm, nosep, itemsep=2pt]
    \item Constant
    \item Variable
    \item Coefficient
    \item Term
    \item Expression
  \end{enumerate}

  \item \small Which part of the expression $2x + 5$ is the constant term?
  \begin{enumerate}[label=(\Alph*), leftmargin=0.6cm, nosep, itemsep=2pt]
    \item $2x$
    \item 5
    \item 2
    \item x
    \item +
  \end{enumerate}

\end{enumerate}
\end{multicols}

\vspace{-0.2cm}
\hrule
\vspace{0.2cm}
\textbf{\small Open Ended Questions (FRQ)}
\begin{enumerate}[label=\textbf{Q\arabic*.}, start=9, itemsep=5pt, topsep=0pt]

  \item \small Draw a diagram to represent the expression $x + 2$ and explain what each part of the expression represents. \vspace{2cm}

  \item \small Explain in your own words the difference between a variable and a constant in an algebraic expression, providing examples of each. \vspace{2cm}

\end{enumerate}

\vfill

% Chef's Tips Footer
\begin{tcolorbox}[colback=blue!5, colframe=mainblue!50, title=\small \textbf{Chef's Tips}, fonttitle=\bfseries, bottom=1mm]
\begin{itemize}[noitemsep, topsep=2pt, leftmargin=0.5cm]
  \item \scriptsize Emphasize the distinction between variables, constants, and coefficients in expressions.\item \scriptsize Use simple, real-world examples to illustrate the concepts of variables and expressions.\item \scriptsize Encourage students to identify and explain the components of basic algebraic expressions.
\end{itemize}
\end{tcolorbox}

\end{document}
  